\documentclass[conference]{IEEEtran}
\usepackage[utf8]{inputenc}
\usepackage[T1]{fontenc}
\usepackage{cite}
\usepackage{booktabs}
\usepackage{graphicx}
\usepackage{url}

\title{Data Understanding v2 for ICU Mortality Dataset as Foundation for Agentic AI in Healthcare}

\author{\IEEEauthorblockN{Hamzah Arman Husni}
\IEEEauthorblockA{Magister Kecerdasan Artifisial\\
Universitas Gadjah Mada\\
Yogyakarta, Indonesia}}

\begin{document}
\maketitle

\begin{abstract}
This paper reports the second refined version (v2) of the data-understanding phase for the ICU Patient Monitoring and Mortality Prediction dataset. The workflow is organized from D1 to D8 in a reproducible sequence, where data card and schema mapping are intentionally positioned at the beginning. The dataset contains 15,000 records and 24 variables, with no missing values and no duplicated rows. Mortality prevalence is 0.2279, indicating the need for imbalance-aware evaluation planning. All figures are generated directly from notebook code with a unified v2 naming convention. The deliverable completion score reaches 100.0/100 based on artifact-level verification.
\end{abstract}

\begin{IEEEkeywords}
Data Understanding, ICU, CRISP-DM, Agentic AI, Healthcare Analytics, Reproducibility
\end{IEEEkeywords}

\section{Introduction}
Data understanding is a critical stage in high-stakes AI applications such as ICU decision support. A technically strong model may still fail if provenance, structure, and quality constraints are not examined in advance.

\section{Dataset and Provenance}
The dataset is sourced from Kaggle and loaded from a single CSV file. Version traceability is ensured using SHA256 checksum. In v2 workflow, the data card is generated at the beginning to align with standard data-understanding principles.

\section{Data Understanding Workflow (D1--D8)}
\subsection{D1: Data Card and Schema}
The v2 workflow documents source, schema, and limitations first. This reduces ambiguity in downstream interpretation.

\subsection{D2: Representativeness and Bias}
Class balance and subgroup-level mortality rates are reported for gender, age group, and admission type.

\subsection{D3: Advanced EDA}
Distribution analysis, missingness inspection, outlier detection, and correlation heatmap are generated as reproducible figures.

\subsection{D4: Semantic and Leakage Validation}
Identifier checks, leakage scan, and pitfall checklist are exported in structured artifact format.

\subsection{D5: Data Quality}
ISO/IEC 25012-inspired quality dimensions are computed and reported.

\subsection{D6: Evaluation Planning}
Split strategy, metric priorities, baseline reference, and leakage guards are explicitly documented.

\subsection{D7: Fairness and Ethics}
Fairness is reported descriptively at data-understanding stage, with predictive fairness deferred to modeling phase.

\subsection{D8: Deliverable Audit}
All artifacts are validated through file-existence checks, resulting in a completion score of 100.0/100.

\section{Results Summary}
Key quantitative observations: dataset size 15,000 $\times$ 24, missing value count 0, duplicate row count 0, mortality prevalence 0.2279, and majority baseline 0.7721.

\section{Conclusion}
The v2 notebook provides a cleaner and fully reproducible data-understanding pipeline, with consistent figure generation and complete deliverable packaging suitable for academic reporting and further modeling phases.

\begin{thebibliography}{00}
\bibitem{b1} M. A. Lones, ``How to avoid machine learning pitfalls,'' arXiv:2108.02497, 2023.
\bibitem{b2} D. B. Rubin, ``Inference and missing data,'' Biometrika, 1976.
\bibitem{b3} J. W. Tukey, \emph{Exploratory Data Analysis}. Addison-Wesley, 1977.
\end{thebibliography}

\end{document}
